
% =============================================================================
% SECTION 6: LIMITATIONS
% =============================================================================
\section{Limitations and Future Work}
\label{sec:limitations}

\owner{Adebayo, with input from all} \status{Done}

\subsection{Model Limitations}
A central simplifying assumption of the model is that firms behave as rational profit-maximizing actors that respond predictably to economic incentives such as permit prices, audit probabilities, and expected penalties. This assumption is standard in economic models of regulatory compliance and enables tractable analysis of equilibrium behavior in the simulated permit market. However, the landscape of advanced AI development includes actors whose motivations may not be reducible to short-term economic payoffs. State-backed research laboratories, geopolitically motivated organizations, and ideologically driven research groups may place substantial weight on strategic, political, or reputational goals. In such contexts, actors might rationally tolerate regulatory penalties or higher compliance costs in pursuit of technological leadership or national advantage. By modeling firms purely as profit maximizers, the framework may therefore underestimate the willingness of some actors to accept regulatory risk or to engage in strategic rule circumvention when non-economic incentives dominate.

The model also makes three simplifying assumptions about how firms decide. First, decisions are myopic: each firm evaluates the deterrence condition period-by-period and does not optimize over future periods, so the framework captures no forward-looking deterrence or strategic anticipation of escalating scrutiny. Second, collateral is costless to post---the upfront deposit $K$ is forfeited only upon a detected violation and otherwise returned in full, with no opportunity cost for the capital it ties up; in practice this financing cost would raise the effective burden of compliance and strengthen deterrence. Third, firms do not update their beliefs: the detection probability is treated as a known, fixed parameter rather than something learned from observed audit outcomes, so the model excludes the belief-based dynamics through which enforcement experience would, in reality, shape future compliance choices. Each of these is a deliberate simplification that preserves the tractability of the one-period deterrence calculus, at the cost of the richer intertemporal behavior a fully dynamic model of firm learning would exhibit.

A further set of limitations concerns the model's institutional scope. The framework assumes a single centralized regulator responsible for implementing and enforcing the compute permit regime. In practice, AI governance is likely to emerge through a patchwork of national regulatory systems, regional agreements, and industry-level standards. Differences in legal frameworks, monitoring capabilities, and enforcement priorities across jurisdictions could lead to inconsistent application of permit rules. Such heterogeneity creates opportunities for regulatory arbitrage, whereby firms shift training workloads to jurisdictions with weaker oversight or lower compliance costs. The current framework abstracts away from these institutional complexities in order to focus on the core dynamics of permit triggers and audit mechanisms. While analytically useful, this simplification limits the model's ability to capture the strategic interactions that arise when multiple regulators compete or cooperate in overseeing a globally distributed AI development ecosystem.

The model also does not incorporate secondary permit markets or mechanisms for permit banking across regulatory periods. In many established regulatory systems such as emissions trading regimes, firms can buy, sell, or accumulate permits over time. These mechanisms introduce additional economic dynamics, including speculative behavior, liquidity constraints, and strategic hoarding of permits in anticipation of future scarcity. Secondary markets may also improve overall efficiency by allowing firms with lower compliance costs to sell unused permits to firms facing higher marginal costs. By excluding these mechanisms, the current model treats permits as static allocations tied to individual firms within each simulation round. This approach clarifies the baseline incentive structure of the permit system but does not capture the richer market dynamics that could arise if permits become tradable financial instruments.

A further limitation concerns the threshold calibration procedure used to determine permit triggers under uncertainty. The calibration approach relies on estimated distributions of firm behavior, monitoring accuracy, and compute usage in order to determine regulatory thresholds that activate additional oversight. While this method improves the systematic design of trigger rules, its effectiveness depends heavily on the accuracy of the assumptions used in the calibration process. In rapidly evolving technological environments such as frontier AI development, the underlying distributions of compute costs, model training scale, and firm behavior may shift over time. If these parameters change substantially, thresholds calibrated using earlier data may no longer reflect the intended balance between regulatory oversight and operational flexibility. As a result, regulators would likely need to periodically update or recalibrate permit thresholds as monitoring technologies, compute markets, and development practices evolve.

A final limitation concerns detection. The framework treats the Type-II error rate (false negatives in violation detection) as a fixed exogenous parameter. In reality, the probability that a violation goes undetected depends on the technological and institutional capabilities available to regulators. Advances in monitoring infrastructure such as improved hardware telemetry, mandatory reporting from cloud service providers, or cryptographic verification of compute usage could substantially reduce the likelihood of undetected violations over time. At the same time, regulated firms may develop increasingly sophisticated techniques for distributing workloads across infrastructure or masking compute usage to evade detection. By modeling detection accuracy as constant, the current framework does not capture this dynamic co-evolution between regulatory monitoring technologies and strategic firm behavior. Future extensions of the model should therefore treat detection accuracy as an endogenous variable that evolves in response to investments in monitoring technology and enforcement capacity.

\subsection{Simulation Limitations}
The simulation environment relies on parameter values calibrated to contemporary estimates of frontier AI development costs, compute availability, and regulatory enforcement capacity. While these parameters allow the model to approximate current conditions in large-scale machine learning development, the underlying technological landscape is changing rapidly. Improvements in hardware efficiency, changes in global semiconductor supply chains, and innovations in training algorithms could significantly alter the cost structure of frontier AI development over relatively short time horizons. As a result, the quantitative outcomes generated by the simulation should be interpreted as illustrative scenarios rather than predictive forecasts. The primary goal of the simulation is to illuminate qualitative relationships between regulatory triggers, monitoring accuracy, and firm behavior, rather than to provide precise estimates of future compliance outcomes.

Another limitation arises from the scale and temporal scope of the simulation. The model simulates a market consisting of 20 firms interacting over 10 steps. While this structure is sufficient to illustrate key dynamics, such as compliance incentives, permit allocation effects, and the interaction between auditing and detection errors, it remains a highly stylized representation of the real-world AI development ecosystem. In practice, the market for large-scale machine learning systems involves a diverse range of actors, including major technology firms, specialized AI startups, academic research groups, and cloud infrastructure providers. These actors operate over longer time horizons and may enter or exit the market as technological opportunities and regulatory conditions evolve. Expanding the simulation to incorporate a larger number of firms and longer time horizons would allow future work to capture richer patterns of strategic adaptation, including firm entry, consolidation, and long-term investment decisions.

A final limitation concerns the interactive dashboard developed to visualize simulation outcomes. The dashboard is intended primarily as a demonstration and exploratory tool for illustrating how the model behaves under different parameter configurations. It enables users to observe the effects of changes in audit rates, permit thresholds, and error probabilities on system-level outcomes such as compliance rates and regulatory costs. However, the current implementation is not designed for operational deployment in real-world regulatory settings. A production-level system would require extensive validation, formal verification of model assumptions, robust security safeguards, and integration with real-time data sources such as compute telemetry or cloud provider reporting systems. In addition, any operational use of such a system would need to account for adversarial behavior from regulated entities, which may attempt to exploit weaknesses in the monitoring or reporting infrastructure.

\subsection{Future Work}
One important direction for future research is the extension of the framework to multi-jurisdiction regulatory environments. AI development is inherently global, with training infrastructure, research teams, and corporate headquarters often distributed across multiple countries. In such settings, firms may strategically allocate training workloads across jurisdictions to minimize regulatory costs or reduce the probability of detection. Modeling these dynamics would require incorporating multiple regulators with potentially different permit thresholds, enforcement capabilities, and monitoring technologies. Such an extension would make it possible to study phenomena such as jurisdiction shopping, regulatory competition, and the potential benefits of international coordination mechanisms designed to harmonize compute governance standards.

Another promising avenue for future work involves incorporating reputation mechanisms and adaptive audit strategies into the regulatory model. In many regulatory contexts, enforcement agencies adjust monitoring intensity based on the past compliance behavior of regulated entities. Firms with a history of violations may face increased scrutiny, while consistently compliant actors may benefit from reduced audit frequency or simplified reporting requirements. Introducing such dynamic enforcement strategies into the model could significantly alter the incentives faced by firms. Reputation-based monitoring could encourage long-term compliance by increasing the expected costs of repeated violations, while also allowing regulators to allocate enforcement resources more efficiently across firms with different risk profiles.

Future research should also explore competitive racing dynamics among firms developing frontier AI systems. The current model treats firms as largely independent decision-makers responding to regulatory incentives, but in reality AI development often occurs within highly competitive technological races. When one firm makes progress toward a significant capability milestone, competitors may respond by increasing investment, accelerating training schedules, or scaling up compute usage in order to maintain technological parity. These racing dynamics could interact with regulatory thresholds in complex ways, potentially creating pressure for firms to operate near the boundary of permitted compute levels. Incorporating competitive dynamics into the model would allow researchers to examine whether permit-trigger systems remain robust under conditions of intense technological competition.

These dynamics also interact closely with the institutional governance mechanisms discussed in Section 5, which outlines potential international coordination frameworks for compute oversight. Integrating such governance proposals into the simulation framework could provide a useful testbed for evaluating how cooperative monitoring arrangements or harmonized permit standards might reduce incentives for jurisdiction shopping.

Finally, the long-term usefulness of the framework will depend on empirical validation using real-world data as compute governance regimes begin to emerge. Governments and international organizations are increasingly exploring policies requiring reporting or registration of large-scale training runs. As such policies are implemented, data may become available regarding compute usage patterns, audit outcomes, and compliance behavior among regulated entities. Integrating these empirical observations into the model would allow researchers to test the assumptions underlying the simulation and refine parameter estimates. Over time, empirical validation could transform the current conceptual model into a more robust decision-support tool capable of informing real-world regulatory design.

 